\section{Canonical normalization and source reflection}\label{sec:normalization}

For completeness we derive the classical normalization used in
\eqref{eq:xi}. Let
\[
 \psi(x)=\sum_{n=1}^\infty e^{-\pi n^2x},\qquad
 \Theta(x)=1+2\psi(x)\quad(x>0).
\]
The Fourier transform of $e^{-\pi x v^2}$, with kernel
$e^{-2\pi\ii v\xi}$, is $x^{-1/2}e^{-\pi\xi^2/x}$.
Poisson summation for this rapidly decaying Gaussian therefore gives
\begin{equation}\label{eq:Jacobitheta}
 \Theta(x)=x^{-1/2}\Theta(1/x),\qquad
 \psi(x)=x^{-1/2}\psi(1/x)+\tfrac12(x^{-1/2}-1).
\end{equation}
Set $f(u)=e^u\psi(e^{4u})$ for real $u$. Direct differentiation of the
locally uniformly convergent series gives
\begin{equation}\label{eq:Phioperator}
 \frac{f''(u)-f(u)}8
 =2e^{9u}\psi''(e^{4u})+3e^{5u}\psi'(e^{4u})=\Phi(u).
\end{equation}
The theta identity gives $f(-u)=f(u)+\sinh u$. Applying
$\partial_u^2-1$, which annihilates $\sinh u$, proves
$\Phi(-u)=\Phi(u)$. Differentiating at zero also gives
$-f'(0)=f'(0)+1$, or $f'(0)=-1/2$.
These identities concern the complete series.

For $\re s>1$, absolute convergence and the gamma integral give
\[
 \pi^{-s/2}\Gamma(s/2)\zeta(s)
 =\int_0^\infty\psi(x)x^{s/2-1}\dd x.
\]
Split at $x=1$, apply \eqref{eq:Jacobitheta} on $(0,1)$, and substitute
$x\mapsto1/x$ in its decaying part. The elementary remaining integral is
$1/(s-1)-1/s=1/(s(s-1))$. Hence
\begin{equation}\label{eq:completedMellin}
 \pi^{-s/2}\Gamma(s/2)\zeta(s)
 =\frac1{s(s-1)}+
 \int_1^\infty\psi(x)
 \left(x^{s/2-1}+x^{(1-s)/2-1}\right)\dd x.
\end{equation}
The last integral is entire in $s$ by its exponential decay. Multiplying
by $s(s-1)/2$ thus continues $\xi$ as an entire function. Put
$v=2s-1$ and $x=e^{4u}$ in \eqref{eq:completedMellin}. This gives
\begin{equation}\label{eq:ximellin}
 \xi(s)=\frac12+(v^2-1)\int_0^\infty f(u)\cosh(vu)\dd u.
\end{equation}
Twice integrating by parts, using the super-exponential decay of $f,f'$
at $+\infty$ and $\sinh(0)=0$, yields
\begin{align*}
 \int_0^\infty(f''-f)\cosh(vu)\dd u
 &=-f'(0)+(v^2-1)\int_0^\infty f(u)\cosh(vu)\dd u\\
 &=\xi(s).
\end{align*}
All these integrals converge for every complex $v$, so analytic
continuation introduces no new restriction. Finally take
$s=(1+\ii z)/2$, for which $v=\ii z$, and use
\eqref{eq:Phioperator}. The result is
\[
 H(0,z)=\int_0^\infty\Phi(u)\cos(zu)\dd u
       =\frac18\xi\bigl((1+\ii z)/2\bigr),
\]
as claimed. This derivation also fixes the factor $1/8$ used by the
energy and numerical records; a scalar change of $H$ would rescale its
reciprocal energy by the inverse scalar modulus.
